% Generated draft: requires mathematical and source review.
\documentclass[11pt]{article}
\usepackage[margin=25mm]{geometry}
\usepackage{fontspec}
\setmainfont{Latin Modern Roman}
\usepackage{amsmath,amssymb,mathtools,unicode-math}
\setmathfont{Latin Modern Math}
\usepackage{enumitem,needspace,xurl,hyperref}
\setlength{\parindent}{0pt}
\setlength{\parskip}{5pt}
\setlength{\emergencystretch}{3em}
\newcommand{\N}{\mathbb{N}}
\newcommand{\Z}{\mathbb{Z}}
\newcommand{\Q}{\mathbb{Q}}
\newcommand{\R}{\mathbb{R}}
\newcommand{\C}{\mathbb{C}}
\newcommand{\F}{\mathbb{F}}
\newcommand{\A}{\mathbb{A}}
\newcommand{\PP}{\mathbb{P}}
\newcommand{\E}{\mathbb{E}}
\newcommand{\eps}{\varepsilon}
\newcommand{\dd}{\,\mathrm{d}}
\begin{document}
\section*{UnsolvedMath ID: 30003280}
\textbf{Source problem code:} catalogue entry
\subsection{Definitions and mathematical statement}
Let \(p_1<p_2<\cdots\) enumerate all positive prime numbers, and let \(q_1<q_2<\cdots\) enumerate all positive primes congruent to \(3\pmod 4\). For each integer \(k\geq 1\), define
\[
P_k=\prod_{i=1}^{k}p_i,\qquad Q_k=\prod_{i=1}^{k}q_i.
\]
For an integer \(n\geq 2\), define its least prime divisor by
\[
\operatorname{lpf}(n)=\min\{\ell:\ell\text{ is prime and }\ell\mid n\}.
\]
The supplied question asks whether one can construct a Euclid-type expression in \(Q_k\) whose value \(E_k\) satisfies
\[
E_k\in\Z_{\geq 2},\qquad \operatorname{lpf}(E_k)=q_{k+1}
\quad\text{for every integer }k\geq 1.
\]
Here \(E_k\) denotes the value of the requested expression; no admissible class of expressions is specified by the source record. In particular, the meanings of ``Euclid-type'' and ``expression in \(Q_k\)'', the permitted auxiliary inputs, and the required uniformity in \(k\) are unresolved. Consequently, the displayed condition records the requested least-prime-divisor property but does not constitute a fully specified existence problem until those restrictions are supplied. The product \(P_k\) is included in the source setup, but no role for it in the requested construction is stated.
\subsection{Short English statement}
The question asks for a Euclid-type construction from the product of the first \(k\) primes congruent to \(3\pmod 4\), with least prime divisor exactly the next such prime. The permitted constructions are unspecified, so the intended existence question requires clarification.
\subsection{Sources}
\begin{itemize}
\item Local UnsolvedMath registry, permanent ID \(30003280\), source problem code \texttt{catalogue entry}: supplied complete source record for ``Euclid-Type Expressions for Primes Congruent to Three Modulo Four'', in the Oberwolfach Reports Open Problems collection. No individual proposer is identified; the supplied proposed year is 2016.
\item Supplied report citation: \emph{Analytic Number Theory (2017)}, Oberwolfach Reports, DOI 10.4171/owr/2016/53, pp. 2975--3030. \url{https://doi.org/10.4171/owr/2016/53}. The record describes Oberwolfach Reports as documenting mathematical workshops and their open-problem discussions.
\item Supplied literature source, ``OWR 53 report'': \url{https://publications.mfo.de/bitstream/handle/mfo/3557/OWR_2016_53.pdf?isAllowed=y&sequence=1}.
\item Literature assessment checked at \texttt{2026-08-21}: the supplied curation labels the problem open and reports that no authoritative later construction or impossibility proof was found. This is supplied dated curation, not an independently verified literature assessment; the record expressly requests independent verification of time-sensitive status and cited claims.
\item The supplied statement status is \texttt{corrected\_verified}, with the qualification that mathematical content was preserved while notation, typography, or omitted definitions were repaired for self-contained publication. This registry label does not establish independent verification of the formulation or mathematical review of this draft.
\end{itemize}
\end{document}
